\documentclass[11pt,a4paper]{article}
\usepackage[T1]{fontenc}
\usepackage{lmodern,amsmath,amssymb,amsthm,microtype,booktabs}
\usepackage[margin=25mm]{geometry}
\usepackage[colorlinks=true,urlcolor=blue,linkcolor=blue]{hyperref}
\newtheorem{theorem}{Theorem}
\newtheorem{corollary}{Corollary}
\newcommand{\E}{\mathcal E}
\newcommand{\dd}{\mathop{}\!\mathrm d}
\title{Transverse masses in the recoil Barnes limit\\and the limits of central localization}
\author{Alan Day \and Chengheng\thanks{AI research assistant (OpenAI).}}
\date{Working research note --- 30 October 2026}
\begin{document}\maketitle
\begin{abstract}
We restore two transverse masses in the internal chiral Barnes graph
obtained after a fixed-mass pure-recoil limit. Their leading effect is
an effective Bessel argument and a modified quadratic phase.
A finite contour decomposition yields a joint estimate uniformly
inside a strict triangle-gap region. At the natural transverse scale,
the two Bessel branches have opposite Gaussian curvatures.
For a specified imaginary-energy regime, even the unnormalized internal
graph grows away from the transverse center. This does not establish
divergence of a physical external integral. The analysis is a
self-reviewed proof draft, with empirical numerical diagnostics.
\end{abstract}
\section{Object and uniform statement}
Use radius one, \(K=\nu^2/(2M)>0\), \(\rho=K^3/\nu^2\),
\(\omega=i\alpha\), and real \(u,v\). The fixed-mass \(\kappa\to\infty\)
limit has already been taken. The masses are
\(m_{1,3}=M+\omega/2\pm u\), \(m_{2,4}=M-\omega/2\pm v\).
The edge ordering is \((m_1,\nu,m_2,m_3,\nu,m_4)\).
The sets in CLP (3.46) \cite{clp} give
\(a_j=M+\nu/2+(\pm u\pm v)/2\) and
\(r_j=M+\nu\pm\omega/2,2M\).
Define
\begin{align}
 w_{\epsilon\eta}&=(\nu-\epsilon u-\eta v)/2,\\
 k&=\frac{2\sqrt{\prod_j w_j}}{M}
 =\frac{\sqrt{[\nu^2-(u+v)^2][\nu^2-(u-v)^2]}}{2M},&
 \lambda&=\frac18\sum_j\frac1{w_j}.
\end{align}
For fixed \(0\le\theta<1\), \(|u|+|v|\le\theta\nu\) implies
\((1-\theta^2)K\le k\le K\), \(w_j\ge(1-\theta)\nu/2\),
and \(\lambda\le[(1-\theta)\nu]^{-1}\).
Let
\[
 J=-e^{i\pi\alpha}kK_\alpha(k)-i\pi kI_\alpha(k),\qquad
 \E=kK_\alpha(k)+\pi kI_\alpha(k),\qquad S_k=k^2+\alpha^2-1.
\]
\begin{theorem}
Fix \(\theta<1\) and \(Y<\infty\). For sufficiently large \(K\) and
sufficiently small \(\rho\), the normalized internal integral defined
below satisfies, uniformly for \(0\le\alpha\le YK\) and
\(|u|+|v|\le\theta\nu\),
\[
 |I-e^{i\lambda S_k}J|\le C_{\theta,Y}
       (\rho+e^{-c_{\theta,Y}K})\E .
\]
The constants are not uniform as \(\theta\to1\).
\end{theorem}
This includes cancellation points of \(J\); relative error requires
\(|J|\ge\epsilon\E\). It neither specifies external states and operator
algebras nor supplies physical insertion normalization or finite-\(G\)
dynamics. No extra vertex is multiplied onto a full \(6j\) expression.

\newpage
\section{Exact kernel, multiplier, and positive tail}
Set \(b=\tfrac12-i(M-\nu)\), \(c=\tfrac12-i(M+\nu)\), and
\begin{align*}
 C(u,v)&=e^{-\pi(M+\nu)}\prod_j\Gamma(-iw_j)\frac{\Gamma(b)}{\Gamma(c)},\\
 Q(t)&=e^{-i\pi t}\prod_j\frac{\Gamma(t-iw_j)}{\Gamma(-iw_j)}
 \frac{\Gamma(b-t)\Gamma(c)}{\Gamma(b)\Gamma(c+t)}
 \Gamma\left(\frac{1+\alpha}2-t\right)
 \Gamma\left(\frac{1-\alpha}2-t\right).
\end{align*}
The contour rises from \(-N-i\infty\) to \(-N+i\), then runs right.
We use the internal analytic continuation specified by this contour,
and \(I=(2\pi i)^{-1}\int Q\,\dd t=\mathfrak B/C\).
This is not a new prescription for the physical external cycle.
Let
\[
 G=(k/2)^{2t}\Gamma\left(\frac{1+\alpha}2-t\right)
                 \Gamma\left(\frac{1-\alpha}2-t\right),\quad
 V=e^{-i\pi t}(k/2)^{2t}
       \frac{\Gamma((1+\alpha)/2-t)}{\Gamma((1+\alpha)/2+t)} .
\]
Then \(Q=Re^{-2\pi it}G\), where
\[
 R=e^{i\pi t}(k/2)^{-2t}\prod_j\frac{\Gamma(t-iw_j)}{\Gamma(-iw_j)}
        \frac{\Gamma(b-t)\Gamma(c)}{\Gamma(b)\Gamma(c+t)} .
\]
Reflection gives the exact, integer-\(\alpha\) regular identity
\[
 e^{-2\pi it}G=-e^{i\pi\alpha}G+2\pi e^{i\pi\alpha/2}V.                 \tag{3}
\]
On \(|t|\le CK\), use logarithms continued from zero shift.
The Gamma expansion \cite{gamma} yields
\[
 \log R=4i\lambda(t^2-t)+O_{C,\theta}(\rho).                         \tag{4}
\]
Indeed, the four numerator remainders are bounded by
\(C\sum_j(|t|+|t|^3)/w_j^2\).
The \(b,c\) powers and shifts contribute at most
\(C|t|(\nu^2/M^2+1/M)+C(|t|+|t|^2)/M=O(\rho)\).
Thus \(|R|\le \exp[O(\rho+K^2/\nu)]=\exp[o(K)]\) on the rectangle.
The leading phase is not expanded in powers of \(\lambda K^2\).

The exact recurrence is
\[
 \frac{Q(t+1)}{Q(t)}
 =\frac{\prod_j(t-iw_j)}
 {[(t+\tfrac12-i\nu)^2+M^2][(t+\tfrac12)^2-\alpha^2/4]} .             \tag{5}
\]
Put \(d_0=\max(\alpha^2/4-1,0)\), \(A=\max(2,8K,2d_0)\),
\(L_0=M-\nu+1\), \(a=\max_j(w_j-1)\le\nu\).
For \(\nu\ge16K\),
\[
 (1+1/(2A))^2(2a^2+a^4/A^2)
 \le\frac{75}{128}\frac{\nu^4}{16K^2}\le L_0^2 .
\]
Factoring the denominators of (5), as in the symmetric estimate,
gives \(|Q(x+1+i)/Q(x+i)|\le[x/(x+1)]^{3/2}\) for \(x\ge A\).
The increased threshold \(8K\), rather than \(2K\), is essential
to this particular bound for distinct \(w_j\).

\newpage
\section{Finite contours and centered branch moments}
Let \(D=\max(8,Y^2/2)\) and \(\nu^2\ge32(D+1)K^3\).
For \(R_0\ge(32+2Y)K\), the transition region
\(R_0\le x\le A+1\), if nonempty, obeys
\[
 |Q(x+1+i)/Q(x+i)|
 \le2[(A+1)/L_0+a^2/(R_0L_0)]^2\le1/8 .
\]
Both terms in brackets are at most \(1/8\).
Combining this with the far tail gives
\[
 \int_{R_0}^{\infty}|Q(x+i)|\,\dd x
 \le10(R_0+1)\sup_{[R_0,R_0+1]}|Q(x+i)|.
\]
Set \(L=\sqrt{k^2+\alpha^2}/2\), \(d=\tfrac12-L\), \(N=\lceil L\rceil\).
For large \(K\), \(N>\alpha/2+1/2\). For every \(w>0\) and \(s\ge0\),
the exact integer Gamma modulus gives
\[
 \left|\frac{\Gamma(-N-i(w+s))}{\Gamma(-N-iw)}\right|
 \le e^{-\pi s/2}.
\]
Applying this separately to all four factors, and the half-integer
product to \(b,c\), gives
\[
 |Q(-N-is)|\le|Q(-N)|e^{-\Lambda s},\quad
 \Lambda=2\pi-N[(M-\nu)^{-1}+(M+\nu)^{-1}].
\]
The real-axis multiplier has leading modulus one, and real Stirling
at the \(G\) saddle gives \(|Q(-N)|\le C_{\theta,Y}\E/\sqrt K\).

Choose \(r\) within one unit of \((32+2Y)K\), at least \(1/8\)
from both real pole lattices. Set \(T=H_{\theta,Y}K\), with fixed large
\(H_{\theta,Y}\). Eventually \(T<\min w_j,M-\nu\).
Move the \(G\) branch to the \emph{finite} line \(d+i[-T,T]\);
close \(V\) clockwise below. Only
\(t_n=(1+\alpha)/2+n<r\) are enclosed, with positive weights
\[
 \omega_n=(k/2)^{1+\alpha+2n}/[n!\Gamma(1+\alpha+n)] .
\]
The resulting formula is
\[
 I=-\frac{e^{i\pi\alpha}}{2\pi}\int_{-T}^{T}R(d+is)G(d+is)\,\dd s
       -2\pi i\sum_{t_n<r}\omega_nR(t_n)+O(e^{-cK}\E).              \tag{6}
\]
For completeness, upper and lower connectors have logarithmic modulus
at most \(-\pi HK+O_{\theta,Y}(K[1+\log(1+H)])\).
The reflected right edge is exponentially small since \(r\gg k,\alpha\)
and avoids both pole lattices. On the lower and right closing edges
\(|V|\le C|G|\). The multiplier is only \(\exp[o(K)]\);
\(\E\ge c_\theta\sqrt K e^{-K}\) absorbs all connectors and both
original tails. No contour crosses the poles at height \(\min w_j\)
or the distant \(b-t\) poles.

Put \(A(t)=4(t^2-t)\).
The Bessel equation at fixed \(\alpha\), with
\(p_n=2\omega_n/[kI_\alpha(k)]\), gives
\[
 \mathbb E_p(A-S_k)=0,\qquad
 \mathbb E_p(A-S_k)^2=8k^2\mathbb E_p t_n=O_{\theta,Y}(K^3).
\]
For the other branch, \(q(s)=A(d+is)-S_k=-4s^2-8iLs\);
its full auxiliary complex moment \(\int qG\,\dd s\) vanishes.
The Gamma product gives Gaussian modulus \(e^{-cs^2/K}\) on the
finite line and \(\int|G|\,\dd s\le CkK_\alpha(k)\).
Taylor remainder costs \(\lambda^2|q|^2e^{8\lambda L|s|}\).
Completing the square costs \(e^{O(\lambda^2L^2K)}=e^{O(\rho)}\).
Both branch errors are therefore \(O(\rho)\) times their own Bessel
magnitudes. Auxiliary tails are exponentially small. Together with
(4)--(6), this proves Theorem 1.

\newpage
\section{Natural transverse scale and normalization}
Take \(u=\sqrt M\,\xi\), \(v=\sqrt M\,\zeta\), \(r_\perp^2=\xi^2+\zeta^2\),
and keep \(\alpha=Ky\) fixed when changing a Bessel argument.
Exactly,
\[
 k^2=K^2-Kr_\perp^2+(\xi^2-\zeta^2)^2/4 .
\]
For bounded \(\xi,\zeta\),
\[
 k=K-\frac{r_\perp^2}{2}-\frac{\xi^2\zeta^2}{2K}+O(K^{-2}).
\]
Define
\[
 \Phi(x,\alpha)=\sqrt{x^2+\alpha^2}
       -\alpha\operatorname{arsinh}(\alpha/x),\qquad
 \partial_x\Phi=\sqrt{1+\alpha^2/x^2}.
\]
The positive-axis Bessel expansions \cite{bessel} have exponent
\(+\Phi\) for \(I_\alpha\) and \(-\Phi\) for \(K_\alpha\).
Consequently
\begin{align}
 \frac{kI_\alpha(k)}{KI_\alpha(K)}
   &\longrightarrow e^{-\sqrt{1+y^2}r_\perp^2/2},\notag\\
 \frac{kK_\alpha(k)}{KK_\alpha(K)}
   &\longrightarrow e^{+\sqrt{1+y^2}r_\perp^2/2}.                  \tag{7}
\end{align}
For bounded \(y\), \(r_\perp\le B\), \(B\ge1\) and
\(B=o(K^{1/4})\), the logarithmic remainders in (7) are
\(O_Y(B^4/K)\). The fixed-order large-argument expansion covers \(y=0\);
the positive real saddle is uniformly nondegenerate on compact \(y\)
intervals. Differentiating along \(\alpha=yx\) would give a different
and incorrect transverse curvature.

Equation (7) first describes auxiliary branch magnitudes.
The finite analytic branches have an additional \(O(\rho)\) error.
The full sum has no uniform relative ratio at cancellation points.
For fixed bounded \(r_\perp\),
\[
 \lambda S_k-\frac{K^2+\alpha^2-1}{\nu}
 =O_Y(Kr_\perp^2/\nu+r_\perp^4/\nu)\longrightarrow0 .
\]
This does not dispose of the separate normalization phase.

The exact identity
\(|\Gamma(-iw)|^2=\pi/[w\sinh(\pi w)]\) and
\(\sum_jw_j=2\nu\) give
\[
 \left|\frac{C(u,v)}{C(0,0)}\right|
 =\frac Kk\,
 \frac{(1-e^{-\pi\nu})^2}
 {\sqrt{\prod_j(1-e^{-2\pi w_j})}} .                              \tag{8}
\]
It tends to one on a bounded natural-scale box. On a fixed
\(\theta\)-region it is \(K/k\) times
\(1+O_\theta(e^{-\pi(1-\theta)\nu})\).
Thus this modulus does not suppress the growing branch.
Its \emph{complex phase} need not tend to zero: we retain the exact
Gamma product, not an expansion of its transverse phase without
additional scale assumptions. For complex masses the two chiral factors
must also be analytically continued separately, not replaced by an
unjustified modulus square.

For \((\xi,\zeta)=(1,1/2)\), \(y=2\), the limiting branch ratios
are \(e^{-5\sqrt5/8}\) and \(e^{+5\sqrt5/8}\).
At \(K=256\), direct Bessel evaluation gives approximately
\(0.246053\) and \(4.04631\), respectively.
The opposite signs are a property of the full two-branch description;
they do not invalidate the earlier result for a specified local saddle.

\newpage
\section{A failure of uniform central decay}
Let \(F(y)=\sqrt{1+y^2}-y\operatorname{arsinh}y\), with positive root
\(y_c\simeq1.50887956\).
Fix \(y>y_c\), \(p=u/\nu\), \(h=v/\nu\), and
\(0<|p|+|h|\le\theta<1\).
Write
\[
 q=\frac kK=\sqrt{[1-(p+h)^2][1-(p-h)^2]}<1,\quad
 \phi_y(q)=\sqrt{q^2+y^2}-y\operatorname{arsinh}(y/q).
\]
Since \(\phi'_y(q)=\sqrt{q^2+y^2}/q>0\),
\(\Delta=\phi_y(1)-\phi_y(q)>0\) and
\(\phi_y(q)<F(y)<0\).
The \(K_\alpha\) branch exponentially dominates at both points.
The envelope bound in Theorem 1 is then a relative bound.

\begin{corollary}
In the above regime, as \(K\to\infty\) and \(\rho\to0\),
the unnormalized internal chiral graph satisfies
\[
 \left|\frac{\mathfrak B(u,v)}{\mathfrak B(0,0)}\right|
 =\left(\frac{1+y^2}{q^2+y^2}\right)^{1/4}
 e^{K\Delta}\left[1+O_{\theta,y}(\rho+K^{-1}+e^{-cK})\right].
                                                                    \tag{9}
\]
\end{corollary}
To see the prefactor, the ratio of \(kK_\alpha(k)\) to
\(KK_\alpha(K)\) contributes
\(q[(1+y^2)/(q^2+y^2)]^{1/4}\).
Equation (8) cancels \(q\), while both leading quadratic phases have
unit modulus. The subleading branch is exponentially small.

Equation (9) excludes a uniform decaying central Gaussian bound on
this real transverse region for this internal object.
It does \emph{not} prove divergence of a physical external integral.
An external cycle, the second chiral factor, physical external legs,
and insertion normalization have not been included.

For \(y=2,p=0.3,h=0.2\),
\(q=0.86168439698\) and \(\Delta=0.32833646092\).
The saved symmetric evaluations and the new transverse evaluations,
with the exact normalization modulus (8), give:
\begin{center}
\begin{tabular}{rrrr}
\toprule
\(K\)&Numerical internal ratio&Leading expression (9)&Relative difference\\
\midrule
8&13.8842&14.0119&0.009114\\
16&192.471&193.756&0.006633\\
32&36871.8&37048.6&0.004771\\
\bottomrule
\end{tabular}
\end{center}
These computations use the finite-branch formula; the theorem controls
omitted connectors asymptotically. The table is not a finite-\(K\)
certified total-error statement.

The immediate research implication is to settle how this branch enters
the physical external transform before integrating a local Gaussian.
The insertion-angle data in CSTY (2.73) \cite{csty} do not by themselves
supply this contour prescription.
The strict-gap theorem also leaves \(|u|+|v|\to\nu\) and beyond open.
Those limits must not be taken inside its constants.

\newpage
\section{Validation, provenance, and next step}
The primary program records 67 checks:
three effective-parameter products; eighteen raw-kernel,
recurrence and exact-split identities; three normalization-modulus
identities; sixteen tail samples; four finite-contour and refinement
checks; twelve joint quadrature refinements; eight precision or extent
checks; and three natural-scale convergence comparisons.
A separate program records sixteen higher-precision multiplier checks.
No failed scientific run was overwritten.

The joint path is \(\nu=8K^{7/4}\), so
\(\rho=1/(64\sqrt K)\to0\) while \(K^2/\nu=K^{1/4}/8\) grows.
We use \(K=8,16,32\), \(y=0,2\), and either
\((u,v)=(\sqrt M,\sqrt M/2)\) or \((0.3\nu,0.2\nu)\).
Quadrature orders are 20 and 32, with 55 and 75 decimal digits.
The largest \(K\) is checked at 95 digits and with height
\(8K\) enlarged to \(10K\), subject to a safe pole-height cap.
The maximum envelope changes are:
\[
 1.755\times10^{-18}\quad\hbox{(quadrature)},\qquad
 3.553\times10^{-70}\quad\hbox{(precision)},\qquad
 2.339\times10^{-45}\quad\hbox{(height)}.
\]
The largest negative-branch and positive-branch condition numbers are
approximately \(1.0804\) and \(1.0263\).
The finite-contour identity retains every connector; its maximum
normalized residual is \(1.011\times10^{-44}\).
Multiplier residuals on \(|t|=O(K)\) are also recorded up to \(K=1024\).
These are empirical checks, not interval enclosures or certified
asymptotic thresholds.

The derivation uses the exact mass sets and contour from CLP,
the standard Gamma and Bessel formulas, and the previously recorded
symmetric contour decomposition. The present work tracks distinct
\(w_j\), changed tail constants, and the effective transverse scale.
It does not establish novelty of these special-function methods.

The next task is the physical external interface: verify the
two chiral continuations, their common external cycle, and the insertion
and disconnected normalizations in the original sources.
Only then decide whether the near-triangle boundary layer is required.
The finite-\(\kappa\) joint limit, finite-\(G\) physical four-point
function, and complete de Sitter holographic dual remain open.

\small
Reproducible records in the project root:
\begin{quote}\footnotesize
code/recoil\_transverse.py\\
code/transverse\_consequences.py\\
results/stage4/recoil-transverse/diagnostics.json\\
results/stage4/recoil-transverse/consequences.json\\
notes/stage4-recoil-transverse.md
\end{quote}
\begin{thebibliography}{9}
\bibitem{clp} Chandra, Liska and Post,
\href{https://arxiv.org/html/2610.30247v1}{arXiv:2610.30247v1},
(3.46), (3.47), (3.58). Local archived version.
\bibitem{csty} Chen, Stanford, Tang and Yang,
\href{https://arxiv.org/html/2607.14042v1}{arXiv:2607.14042v1},
(2.73) and the previously recorded external-transform dictionary.
\bibitem{gamma} NIST DLMF,
\href{https://dlmf.nist.gov/5.4}{5.4},
\href{https://dlmf.nist.gov/5.8}{5.8},
\href{https://dlmf.nist.gov/5.11}{5.11}.
\bibitem{bessel} NIST DLMF,
\href{https://dlmf.nist.gov/10.25}{10.25},
\href{https://dlmf.nist.gov/10.40}{10.40},
\href{https://dlmf.nist.gov/10.41}{10.41}.
\end{thebibliography}
\end{document}
